% Appendix A: short answers to the hand-calculation exercises, numbers only. The reasoning, and
% the plausible wrong answers, stay in the worked solutions in the companion repository. Add each
% new chapter's hand-calculation exercises here as it is typeset.

\chapter{Short Answers}
\label{app:answers}

{\small
\noindent The numbers only, for the hand-calculation exercises, so that an answer worked on paper
can be checked on paper. How each one is reached, and the plausible wrong answers, will be in the
worked solutions, which are published later in the companion repository.

\newcommand{\answer}[2]{\par\addvspace{1.2ex}\needspace{3\baselineskip}\noindent
  {\bfseries\color{accent}Exercise~\ref{#1}}\quad{\bfseries #2}\par\nopagebreak\smallskip\noindent}
\newcommand{\pt}[1]{\textbf{#1)}~}

\answer{c1:ex:assemble}{Assembling four instructions by hand}
\pt{a}\code{0xE20A}\quad \pt{b}\code{0xE28A}\quad \pt{c}\code{0xEFFF}\quad
\pt{d}\code{0xE043}\quad \pt{e}Rejected: \code{ldi} can only name \code{r16} to \code{r31}. A
word forced out anyway would load \code{r31}.

\answer{c1:ex:counting}{Counting by hand}
\pt{a}10 cycles\quad \pt{b}52 cycles\quad \pt{c}$6n + 10$\quad \pt{d}\code{cp} 4 times,
\code{breq} 4 times, 1 of them taken\quad \pt{e}$7n + 10$: the extra \code{inc} on each trip is
the extra cycle per bit.

\answer{c2:ex:translate}{Translating pin numbers}
\pt{a}\begin{tabular}[t]{@{}lllllll@{}}
  Pin & Port & Bit & Mask & \code{PINx} & \code{DDRx} & \code{PORTx} \\
  3 & D & 3 & \code{0x08} & \code{0x0029} & \code{0x002A} & \code{0x002B} \\
  8 & B & 0 & \code{0x01} & \code{0x0023} & \code{0x0024} & \code{0x0025} \\
  13 & B & 5 & \code{0x20} & \code{0x0023} & \code{0x0024} & \code{0x0025} \\
\end{tabular}\par\smallskip
\pt{b}\code{0x01}; without the subtraction \code{0x00}, and the LED never lights\quad
\pt{c}\code{PB6} and \code{PB7}, the crystal's pins.

\answer{c2:ex:predict}{Predicting the driver's cost}
\pt{a}30 cycles more\quad \pt{b}35\quad \pt{c}65\quad \pt{d}\code{led_toggle}: one
\code{st} after the shift, with nothing read\quad \pt{e}Only the shift loop's trip count depends on the pin.

\answer{c3:ex:response}{Addresses and response times}
\pt{a}\begin{tabular}[t]{@{}llll@{}}
  Vector & Number & Word & Byte \\
  PCINT0 & 3 & \code{0x0006} & \code{0x000C} \\
  PCINT2 & 5 & \code{0x000A} & \code{0x0014} \\
  WDT & 6 & \code{0x000C} & \code{0x0018} \\
\end{tabular}\par\smallskip
\pt{b}6 cycles, $4 + 2$\quad \pt{c}9, $3 + 4 + 2$\quad \pt{d}7 and 10; \code{jmp} reaches
all of flash and \code{rjmp} only $\pm$2K words\quad \pt{e}110, $4 + 2 + 104$.

\answer{c3:ex:counter}{The shared counter}
\pt{a}\code{0x01FF}, which is 511; the counter held 255, then 256\quad \pt{b}2 cycles; about
one read in eight thousand, so a wrong value every eight seconds or so\quad \pt{c}\code{cli}
before and \code{sei} after, 2 cycles\quad \pt{d}It enables interrupts whatever they were
before; save \code{SREG}, \code{cli}, load, restore \code{SREG}: 3 cycles\quad \pt{e}No: one
\code{lds} is one instruction.

\answer{c4:ex:fourways}{Four ways to read a byte}
\pt{a}\code{0x33}, \code{Z} stays \code{0x0204}\quad \pt{b}\code{0x33}, \code{Z} becomes
\code{0x0205}\quad \pt{c}\code{0x22}, \code{Z} becomes \code{0x0203}\quad \pt{d}\code{0x66},
\code{Z} stays \code{0x0204}\quad \pt{e}(a) and (d)\quad \pt{f}Two cycles each, so the choice is
only about where the pointer is left.

\answer{c4:ex:arrays}{Array arithmetic}
\pt{a}\code{0x0300}, \code{0x0307}, \code{0x030E}, \code{0x0323}\quad \pt{b}7, \code{LED_SIZE}
in \file{led.inc}\quad \pt{c}\code{0x030C}, the sixth byte of entry 1, so the driver reads
\code{0x0100} (\code{0x0n00} for bit $n$) as a port register\quad \pt{d}Every address is legal,
and entry 0 is at the base under either stride\quad \pt{e}\code{0x031E}: the stride is 10\quad \pt{f}292 LEDs, 204
buttons; the real limit is the 14 pins.

\answer{c4:ex:depth}{How deep does it go}
\pt{a}4 bytes\quad \pt{b}15, $0 + 2 + 9 + 2 \times 2$\quad \pt{c}\code{SP} ends at
\code{0x08F0}; the lowest byte used is \code{0x08F1}\quad \pt{d}2 bytes, the program
counter\quad \pt{e}21, $2 \times 3 + 2 + 9 + 2 \times 2$\quad \pt{f}No: it occupies
\code{0x08F0} to \code{0x08F6}, and the stack reaches down to \code{0x08F1}.

\answer{c5:ex:compare}{Prescaler and compare}
\pt{a--c}\begin{tabular}[t]{@{}lll@{}}
  Wanted & Timer1, 16-bit & Timer0, 8-bit \\
  500\,Hz & \code{/1}, \code{OCR = 31999}, exact & \code{/256}, \code{OCR = 124}, exact \\
  440\,Hz & \code{/1}, \code{OCR = 36363}, $-0.0010\%$ & \code{/256}, \code{OCR = 141}, $+0.032\%$ \\
  3\,kHz & \code{/1}, \code{OCR = 5332}, $+0.0063\%$ & \code{/64}, \code{OCR = 82}, $+0.40\%$ \\
\end{tabular}\par\smallskip
\pt{d}500\,Hz: 32000 ticks divide the clock exactly, and the prescalers are powers of
two\quad \pt{e}3\,kHz, 64 times worse: 83 ticks against 5333.

\answer{c5:ex:reach}{What cannot be reached}
\pt{a}$f_{\mathrm{clk}} / (1024 \times \text{counter values})$: 0.2384\,Hz for 16 bits, 61.04\,Hz
for 8\quad \pt{b}Yes: \code{/1024}, \code{OCR1A = 31249}, exact\quad \pt{c}No: 156250 ticks, and
the counter holds 65536\quad \pt{d}The counter runs out of values: a full 8-bit count at
\code{/1024} spans only 1/61 of a second\quad \pt{e}A 1\,kHz interrupt and a software target of
10000\quad \pt{f}16\,MHz, one tick at a prescaler of 1; nothing is shorter than one tick.

\answer{c6:ex:bytes}{Control bytes}
\pt{a}\code{0x0E}, selector 6 plus \code{WDE}\quad \pt{b}\code{0x61}, selector 9 as \code{0x21}
plus \code{WDIE}\quad \pt{c}\code{0x4C}, selector 4 plus both\quad \pt{d}System reset at 16\,ms:
the device resets, starts up, and resets again, for ever\quad \pt{e}A short timeout is a wrong
setting; \code{0x08} is a device that no longer runs.
\par}
